\textbf{Step 1 --- risk parity.} At each month-end $t$ we estimate the covariance matrix $\Sigma_t$ of the six ETFs from
the last 252 daily returns and choose weights $w$ so that every asset class contributes the same share of portfolio risk
\citep{maillard2010}:
\[
RC_i=\frac{w_i(\Sigma_t w)_i}{w^\top\Sigma_t w}=\frac{1}{6}\ \ \text{for all } i,\qquad w_i>0,\quad \textstyle\sum_i w_i=1 .
\]
There is no closed form, so the six equations are solved with SciPy's \texttt{fsolve}, following \citet{spinu2013}. In all
months the risk contributions are equal to within $10^{-8}$; the inverse-volatility fallback was never needed. Volatile
assets (equities, commodities) get small weights and stable ones (Treasuries) large weights, so no asset class dominates the
risk, as equities do in 60/40 \citep{asness2012}.

\textbf{Step 2 --- trend filter.} An asset class is held only if its month-end price is above its 10-month moving average
\citep{faber2007}; otherwise its weight goes to T-bills. Prices tend to trend over months \citep{moskowitz2012}, so the
filter exits asset classes in prolonged declines such as 2008 or the 2022 bond sell-off.

\textbf{Timing and rebalancing.} Weights are set with closing prices on the last trading day of month $t$ and earn the
returns of month $t+1$; no later information is used. Between monthly rebalancing dates holdings drift with prices.
\textbf{Motivation:} in our sample risk parity alone cut the drawdown of an equal-weighted portfolio from \res{ewnt_max_dd}
to \res{ercnt_max_dd}; the trend filter cut it further to \res{strat_max_dd} (Table~\ref{tab:performance}).
